\chapter{重子生成、宇宙的平滑性与哈勃摩擦}

\lectureoverview{两个小的无量纲数组织了今晚的内容。第一个是幸存的正物质--反物质超出，在课程中约化为\(10^8\)分之一。其动力学解释需要重子数破坏、粒子--反粒子不对称以及偏离热平衡：即三个萨哈罗夫条件。第二个是原初密度反差，在解耦时约为\(10^{-5}\)量级。引力会放大而非抹去这类涟漪，因此它们的微小性促使我们引入暴胀。随后我们为暴胀的机制做准备，推导膨胀宇宙中均匀标量场的运动方程，并确认\(3H\dot\phi\)为哈勃摩擦。}

\section{被颠倒过来的熵问题}

今晚我们将开始讲暴胀，但首先必须完成重子生成：正物质相对于反物质的超出的产生。这一框架比任何详细的数值模型都更为可靠。这就是为什么尽管其核心数字尚未从第一性原理推导出来，这一主题仍能被经济地概括。

进入这一问题的历史路径始于1980年，当时Bob Wagoner问Savas Dimopoulos和Leonard Susskind，为什么宇宙的熵如此之大。用宇宙学的语言来说，问题就是：为什么每一个幸存质子对应着如此多的光子。

在极高温下，光子、夸克、反夸克、电子和正电子都大量存在并接近热平衡。夸克和反夸克各自的数密度与光子密度处于同一大致量级。这一观察已经提示我们把问题反过来问：为什么湮灭会留下任何未配对的正物质？

这一结构性论证的优先权属于Andrei Sakharov。Dimopoulos和Susskind在1980年独立得出了密切相关的条件，随后在准备一篇历史性贡献时发现了Sakharov1967年的论文。这些条件如今被公正地以Sakharov的名字命名。

与熵的联系现在可以明确建立。对于黑体辐射，熵与光子数成正比，相差一个量级为1的固定因子。确实，
\begin{equation}
s_\gamma=\frac{4\rho_\gamma}{3T},
\qquad
n_\gamma\propto T^3,
\end{equation}
因此共动热光子气体有\(S_\gamma\propto N_\gamma\)。数值常数不是问题所在；谜题在于巨大的光子与质子之比。

\begin{classroomqa}
\audiencequestion{光子丰度与熵有何关系？}
\lecturerresponse{黑体辐射是热性的，其熵与其平均光子数成正比。在此语境下，问宇宙熵为何如此之大，本质上就是问宇宙微波背景中每个质子为何对应如此多的光子。这一正比关系包含数值常数，但它们与量级谜题无关。\lecturetimestamp{00:04:49--00:05:55}}
\end{classroomqa}

随着宇宙冷却，粒子和反粒子发生湮灭，而辐射则留存下来并最终解耦。因此，令人惊讶的事实并不是光子幸存了下来；而是湮灭留下了一小群未配对的正物质。

因此有用的无量纲数是
\begin{equation}
\eta_B
\equiv
\frac{N_B-N_{\bar B}}{N_\gamma}
\sim
\frac{N_p-N_{\bar p}}{N_p+N_{\bar p}}
\sim 10^{-8},
\label{eq:c08-baryon-asymmetry}
\end{equation}
其中最后一个值是课程中使用的取整量级，而非精确测定。一旦出现一个看起来奇怪的数字，``它从一开始就是这样''就不再是一个令人满意的终点。一个定量机制应当解释为什么答案很小，以及为什么它具有确定的符号。

把超出称为\emph{正物质}而非\emph{反物质}，部分是约定俗成的：正物质是我们由之构成的配对成员。物理陈述仅仅是：在整个可观测区域中，一种符号占主导。

\section{重子数与第一个条件}

重子数的归一化使得质子或中子带有\(B=1\)。用夸克的语言来说，
\begin{equation}
B=\frac{N_q-N_{\bar q}}{3}.
\label{eq:c08-baryon-number}
\end{equation}
因子\(1/3\)是本质性的，因为每个普通重子包含三个价夸克。其他强子也可能携带重子数，但不稳定的强子最终会衰变；净夸克超出是更一般的记账工具。

\begin{classroomqa}
\classquestion{净夸克数应当乘以还是除以三？}
\lecturerresponse{必须除以三。质子有三个夸克且重子数为一，所以\(3B=N_q-N_{\bar q}\)。黑板上的更正是定义的一部分，而不是物理论证的改变。\lecturetimestamp{00:09:23--00:10:17}}
\end{classroomqa}

重子数仅在作为一种可能的守恒量这一点上类似于电荷。电荷还会作为长程库仑场的源；重子数则不会。质子产生电场是因为它带有电荷，而中子具有重子数却没有净电荷。因此，守恒与力的产生在逻辑上是彼此独立的性质。

\begin{classroomqa}
\audiencequestion{重子数与电荷之间的相关区别是什么？}
\lecturerresponse{两者都可以充当记账用的荷，但只有电荷已知会作为长程电磁场的源。重子数不会产生相应的库仑力。\lecturetimestamp{00:11:03--00:12:15}}
\end{classroomqa}

如果\(B\)严格守恒，那么初始的过剩将永远冻结不变。粒子对的产生与湮灭可以分别改变\(N_B\)和\(N_{\bar B}\)，但只能以相等的量改变，从而使二者之差保持不变。因此，从平衡态出发的动力学理论无法成立，除非重子数能够改变。

最简洁的示意性例子是质子衰变。电荷与能量允许我们设想
\begin{equation}
p\longrightarrow e^+ +\gamma,
\qquad
\bar p\longrightarrow e^-+\gamma.
\label{eq:c08-decay-pair}
\end{equation}
这些通道只是重子数破坏的记账模型，而不是对主导实验通道的预言。它们使\(B=+1\)或\(-1\)变为零，同时保持电荷守恒。

\begin{figure}[htbp]
\centering
\begin{coursefitfigure}
\begin{tikzpicture}[x=1cm,y=1cm,>=Latex]
\draw[->,thick] (0,1.7)--(2.0,1.7);
\draw[->,thick] (2.0,1.7)--(4.1,2.55);
\draw[->,thick] (2.0,1.7)--(4.1,0.85);
\node[above] at (0.9,1.7) {\(p,\ B=1\)};
\node[right] at (4.1,2.55) {\(\gamma\)};
\node[right] at (4.1,0.85) {\(e^+\)};
\draw[->,thick] (6.0,1.7)--(8.0,1.7);
\draw[->,thick] (8.0,1.7)--(10.1,2.55);
\draw[->,thick] (8.0,1.7)--(10.1,0.85);
\node[above] at (6.9,1.7) {\(\bar p,\ B=-1\)};
\node[right] at (10.1,2.55) {\(\gamma\)};
\node[right] at (10.1,0.85) {\(e^-\)};
\node at (2.0,0.25) {\(B:1\to0\)};
\node at (8.0,0.25) {\(B:-1\to0\)};
\end{tikzpicture}
\end{coursefitfigure}
\caption{电荷共轭的示意性通道。两种过程都破坏重子数，但若速率相等，则不会产生任何优先的符号。}
\label{fig:c08-decay-redraw}
\end{figure}

\begin{classroomqa}
\audiencequestion{质子能否改为衰变为正电子和中微子？}
\lecturerresponse{这个双费米子末态与单费米子初态的费米子宇称不匹配。示意性的\(e^++\gamma\)末态包含一个费米子，足以用于该对称性论证。现实的质子衰变理论可能包含其他粒子和通道，但当前目的仅在于展示\(\Delta B\ne0\)。\lecturetimestamp{00:14:39--00:15:07}}
\end{classroomqa}

\begin{classroomqa}
\audiencequestion{严格的重子数守恒是否对应于一种对称性？}
\lecturerresponse{是的。守恒量可以与连续的相位对称性相关联，尽管仅仅命名该对称性并不能独立地证明自然界遵守它。可操作的检验仍然是重子数改变过程是否发生，因此，如果仅从守恒定律推断出对称性，论证就变成循环论证。\lecturetimestamp{00:15:08--00:15:46}}
\end{classroomqa}

\begin{classroomqa}
\audiencequestion{质子--反质子湮灭本身能否被视为质子衰变？}
\lecturerresponse{术语是次要的；守恒定律则不是。在\(p+\bar p\to\text{photons}\)中，初始重子数为\(1-1=0\)，末态重子数也为零。湮灭移除了两个粒子而不改变净\(B\)，而方程~\eqref{eq:c08-decay-pair}中的单粒子通道则改变\(B\)。\lecturetimestamp{00:15:48--00:16:47}}
\end{classroomqa}

当今的质子极其稳定：讲义引用的界限约为\(10^{32}\)--\(10^{34}\)年，远远长于宇宙年龄。这一观测并不能确立严格守恒；它只告诉我们，在当今能量下，重子数改变过程极其缓慢。

\section{离散对称性与对偏向的需要}

下一个问题是，定律是否区分物质与反物质。三种离散变换组织了这一答案：
\begin{align}
C&:\ \text{particle}\longleftrightarrow\text{antiparticle},\\
P&:\ \mathbf x\longmapsto-\mathbf x,\\
T&:\ t\longmapsto-t.
\end{align}
这里\(P\)反转空间手性，而\(T\)反转微观影片。宏观第二定律不必看起来具有时间反演不变性，但问题涉及的是底层微观动力学。

\begin{figure}[htbp]
\centering
\includegraphics[width=0.82\textwidth]{lecture_08_figure_03.png}
\caption{黑板上将修正后的定义\(B=(N_q-N_{\bar q})/3\)带入\(C\)即电荷共轭的引入。charge一词之后未完成的书写不作为独立证据使用。}
\label{fig:c08-charge-board}
\lectureframe{00:20:23.280}
\end{figure}

相对论性局域量子场论并不要求\(C\)、\(P\)或\(T\)分别成为对称性。在其标准假设下，它确实要求组合变换\(CPT\)。如果一个过程是可能的，那么将每个粒子替换为其反粒子、反射空间构型并反转运动，就给出另一个允许的过程。

量子力学中的时间反演比仅仅改变\(t\)的符号更为微妙。对于时间反演不变势中的无自旋波函数，
\begin{equation}
i\hbar\frac{\partial\psi}{\partial t}=H\psi,
\qquad
T:\ \psi(\mathbf x,t)\longmapsto\psi^*(\mathbf x,-t).
\end{equation}
改变\(t\to-t\)会反转时间导数；复共轭也会使\(i\to-i\)，从而恢复方程。对于有自旋的粒子，反幺正的时间反演算符带有额外结构，但复共轭在这个初等情形中已经可见。

\begin{classroomqa}
\audiencequestion{对于波函数，时间反演是否既需要复共轭，也需要\(t\to-t\)？}
\lecturerresponse{是的。薛定谔方程说明了原因：仅反转时间导数会使等式一侧不一致地改变，而复共轭还会反转\(i\)的符号。\lecturetimestamp{00:23:58--00:25:04}}
\end{classroomqa}

现在我们可以陈述第一个萨哈罗夫条件。如果宇宙开始时粒子和反粒子的数量相等，那么除非某种过程能够改变\(B\)，否则粒子对的产生和湮灭将保持净重子数为零：
\begin{equation}
\boxed{\text{baryon-number violation}.}
\end{equation}
没有它，对称的初始状态永远无法获得非零的重子过剩。

质子目前的稳定性与这一条件相容。在大统一理论推理中，重子数破坏可以由质量\(M_X\)极重的粒子传递。在远低于\(M_X\)的能量下，有效振幅按逆幂被示意性地压低，
\begin{equation}
\begin{aligned}
\mathcal A_{\Delta B\ne0}&\propto \frac{E^2}{M_X^2},\\
\Gamma_p&\sim \frac{m_p^5}{M_X^4}.
\end{aligned}
\label{eq:c08-heavy-suppression}
\end{equation}
衰变率的估计具有通常的六维标度，只差耦合常数和依赖于通道的因子。非常大的\(M_X\)可以使冷质子实际上永久稳定，同时允许在能量足够高的早期等离子体中发生改变重子数的反应。

\begin{classroomqa}
\audiencequestion{额外的碰撞能量是否已经包含在重质量\(M_X\)中？}
\lecturerresponse{不。\(M_X\)是传递粒子的静质量标度；\(E\)是碰撞中可用的能量。当\(E\ll M_X\)时，按\(E/M_X\)展开的低能展开被强烈压低，而随着热碰撞标度接近重粒子阈值，这种压低会减弱。黑板上启发式写出的字面分母应在这种有效场论意义下理解。\lecturetimestamp{00:32:32--00:34:02}}
\end{classroomqa}

假设初始状态包含数量相等的粒子和反粒子。仅重子数破坏仍然不会产生偏向。每个衰变\(p\to e^++\gamma\)都会有一个电荷共轭伙伴\(\bar p\to e^-+\gamma\)，而且热等离子体也会驱动逆反应。相等的微观概率会产生通过重子数的流量，但没有平均漂移。

\begin{classroomqa}
\audiencequestion{反质子图中的重子数变化是多少？}
\lecturerresponse{入射反质子具有\(B=-1\)，而电子和光子具有\(B=0\)，因此该通道将重子数从\(-1\)变为零。它的共轭通道将\(+1\)变为零。两个通道出现的机会相等，因此没有净偏好。\lecturetimestamp{00:37:24--00:37:39}}
\end{classroomqa}

随机涨落远太小，无法解释观测到的过剩。如果\(N\)个独立事件对两种结果具有相等概率，那么
\begin{equation}
\Delta N\sim\sqrt N,
\qquad
\frac{\Delta N}{N}\sim N^{-1/2}.
\end{equation}
使用讲座中对可观测宇宙重子数的粗略计数\(N\sim10^{80}\)，得到
\begin{equation}
\frac{\Delta N}{N}\sim10^{-40},
\end{equation}
比\(10^{-8}\)低三十二个数量级。抛硬币式的涨落也没有给出任何理由说明为什么同一符号会主导每一个观测区域。

\begin{classroomqa}
\audiencequestion{一个初始时大得多的粒子数，也许是在后来成为我们可观测宇宙的共动区域中的\(10^{160}\)对粒子，能否涨落\(10^{80}\)？}
\lecturerresponse{平方根算术允许这种人为构造的计数，但它不能解释观测到的空间相干性。不受限制的统计涨落自然会产生物质主导和反物质主导的斑块，而不是在整个可见宇宙中只有一种符号。\lecturetimestamp{00:39:38--00:40:18}}
\end{classroomqa}

讲座时代的观测论证使用两类证据。普通原子核，包括氦核，出现在宇宙射线中，而在当时讨论的数据中并未观测到反原子核。仅凭反质子不那么具有决定性，因为高能碰撞可以产生次级反质子。此外，物质星系和反物质星系之间的边界会产生显著的湮灭辐射。所要求的边界信号缺失，强烈不利于可观测区域内星系和反星系等量混合。这些是讲座历史背景下的观测论证，而不是关于我们视界之外区域的断言。

\begin{classroomqa}
\audiencequestion{氦核本身是在宇宙射线中测量的吗？}
\lecturerresponse{是的。带电氦核是直接测量的。详细的宇宙射线统计是一个单独的专门领域，但普通原子核常见与没有相当数量的反原子核群体之间的对比，正是这一论证的要点。\lecturetimestamp{00:41:53--00:42:22}}
\end{classroomqa}

因此我们需要一种微观偏置：\(C\)的破坏，以及在现实的宇宙重子生成讨论中，\(CP\)的破坏。中性\(B\)介子系统提供了一个实验室实例。中性介子及其反粒子可以衰变到共轭末态，且分宽度不相等，
\begin{equation}
\Gamma(B^0\to f)\ne\Gamma(\bar B^0\to\bar f).
\label{eq:c08-cp-partial}
\end{equation}
确切的通道和数值不对称性在黑板上并未固定；可靠的教训是弱衰变中粒子--反粒子不对称性在实验上已被确立。

尽管如此，\(CPT\)定理要求真空中粒子与其反粒子具有相等的质量和相等的总衰变宽度：
\begin{equation}
\Gamma_{\rm tot}(p)=\Gamma_{\rm tot}(\bar p).
\end{equation}
这并不强制每个分支比都一致。一对共轭通道中的差异可以由其他通道中的差异来补偿，例如示意性的电子和{\proseunicodefont μ}子通道，而总和保持相等。

\begin{classroomqa}
\audiencequestion{如果质子和反质子的总衰变率相等，偏置如何产生？}
\lecturerresponse{总宽度的相等并不意味着逐通道相等。此外，重子生成发生在等离子体中：热占据因子、逆反应和散射通道都进入动理学方程。决定性的陈述并不是介质简单地规避了\(CPT\)，而是一旦细致平衡丧失，\(CP\)不对称的部分过程就能产生源项。\lecturetimestamp{00:49:39--00:51:48; 00:53:22--00:57:14}}
\end{classroomqa}

这就是第二个萨哈罗夫条件：
\begin{equation}
\boxed{C\ \text{and}\ CP\ \text{violation}.}
\end{equation}
它为改变重子数的反应提供了方向性。在萨哈罗夫论文发表前不久发现的\(CP\)破坏使这一条件在物理上变得具体。

最小标准模型需要仔细区分。它不允许通常的微扰质子衰变，但其电弱反常允许\(B+L\)的非微扰破坏；在高温下，sphaleron跃迁可以使这种破坏变得高效。这就是黑板上后来自我修正的正确内容。仅标准模型源能否定量地重现宇宙不对称性，则是一个不同且困难得多的问题。

\section{平衡与第三个条件}

还剩下一个障碍。在完全热平衡中，每个反应都与其逆反应相平衡。如果重子数可以改变，其平衡化学势就不能保持任意的初始\(B\)，而细致平衡会抹去任何产生的过量。讲座通过\(CPT\)表达了同一点：静态平衡浴没有宏观时间箭头，因此\(T\)因子无法提供过程与其逆过程之间缺失的区别。

实际宇宙正在膨胀和冷却。因此宏观状态区分了正向和反向时间。只有当膨胀成功地与微观弛豫竞争时，这一点才重要：
\begin{equation}
\Gamma_{\rm reaction}\lesssim H.
\label{eq:c08-freezeout-condition}
\end{equation}
如果\(\Gamma_{\rm reaction}\gg H\)，碰撞会迅速恢复平衡，尽管存在膨胀。如果反应速率降到哈勃速率以下，等离子体就无法重新调整，丰度便冻结出来。

\begin{classroomqa}
\audiencequestion{为什么仅靠膨胀不足以自动使系统脱离平衡？}
\lecturerresponse{一个膨胀非常缓慢的浴会保持任意接近平衡，因为微观反应不断地重新调整它。只有当膨胀或冷却时间\(H^{-1}\)与相关反应时间\(\Gamma^{-1}\)相比很短时，偏离才变得有效。那时，宏观状态的影片才揭示哪个方向是正向。\lecturetimestamp{01:01:22--01:02:26}}
\end{classroomqa}

现在可以汇总三个萨哈罗夫条件：
\begin{enumerate}
\item 重子数破坏，因此\(B\)能够改变；
\item \(C\)和\(CP\)破坏，因此这种改变具有优先符号；
\item 偏离热平衡，因此细致平衡无法抵消这种偏好。
\end{enumerate}
对于所讨论的机制，它们是必要的，并提供了一条结构上足以产生过量的途径。若不借助具体的高能模型、反应网络和宇宙学历史，它们无法确定其大小。

\begin{classroomqa}
\audiencequestion{不对称性的测量值是否约束高能理论？}
\lecturerresponse{原则上是的。实际上，结果取决于许多未知的质量、耦合、相位、反应速率以及早期宇宙学的特征。因此，若不承诺一个详细模型，单一测量数值很难转化为一个狭窄的允许区域。\lecturetimestamp{01:04:20--01:04:52}}
\end{classroomqa}

\begin{classroomqa}
\audiencequestion{重子不对称性与电子不对称性是如何联系起来的？}
\lecturerresponse{电中性将它们联系起来。在这个示意性通道中，移除一个质子会产生一个正电子；反转重子不平衡会使轻子不平衡朝补偿方向移动。更一般地说，重子生成与轻子生成通过电荷守恒和标准模型反应耦合在一起，尽管守恒的组合取决于具体机制。\lecturetimestamp{01:04:54--01:05:59}}
\end{classroomqa}

\begin{classroomqa}
\audiencequestion{观测到的标准模型\(CP\)破坏是否太小，不足以解释重子过剩？}
\lecturerresponse{这是详细的电弱重子生成分析的标准结论，但讲座强调了该计算对模型的依赖性以及高能区众多未知因素。此处可以站得住脚的陈述是：已知的低能\(CP\)破坏确实存在，而一个成功的数值重子生成模型通常需要额外的要素。\lecturetimestamp{01:06:01--01:07:50}}
\end{classroomqa}

\begin{classroomqa}
\audiencequestion{今天的电子是否仅仅来自反质子的湮灭？}
\lecturerresponse{这些示意性通道使这种记账方式颇具诱惑力，但早期等离子体中也包含独立的、近乎对称的轻子群体。人们可以将相关机制表述为重子生成或轻子生成，随后通过反应在两个部门之间转移不对称性。历史上的\emph{重子生成}一词不应被误认为是每个电子唯一的微观谱系。\lecturetimestamp{01:07:51--01:08:44}}
\end{classroomqa}

\begin{classroomqa}
\audiencequestion{为什么讨论通常以重子而非所有轻子来表述？}
\lecturerresponse{重子物质通过原子光谱和其他天文观测来清点，而宇宙中微子部门则远更难直接计数。这一命名既反映了观测上可及的内容，也反映了该学科的历史。\lecturetimestamp{01:08:45--01:09:12}}
\end{classroomqa}

\section{湮灭留下过剩}

在高温下，粒子对的产生与湮灭双向进行，
\begin{equation}
\gamma+\gamma\longleftrightarrow p+\bar p,
\qquad
E_{\gamma\gamma}\gtrsim2m_pc^2\approx2\,\mathrm{GeV}.
\end{equation}
随着热浴冷却，能量足以产生重子对的光子变得指数级稀少。湮灭在能量上仍然允许，但逆产生过程已无法维持平衡。几乎每个质子都会找到一个反质子，而方程~\eqref{eq:c08-baryon-asymmetry}中的微小过剩就是剩下的东西。高能湮灭光子迅速将其能量与热浴分享。

\begin{classroomqa}
\audiencequestion{比值\(10^{-8}\)是当前的值吗？大规模湮灭发生在什么时候？}
\lecturerresponse{这是本课程给出的当前重子与光子数量级之比。更早的时候，重子和反重子都数量众多。一旦逆粒子对产生过程变得无效，它们的对称分量便湮灭，留下净过剩。\lecturetimestamp{01:09:13--01:11:44}}
\end{classroomqa}

\begin{classroomqa}
\audiencequestion{质子--反质子湮灭何时停止？}
\lecturerresponse{当反质子被消耗到可忽略的丰度时，它就停止了。讲座只给出了大致的时间定位：在GeV量级的温度下，远在核合成之前，当然也远在复合之前。讨论中并未推断出精确的时间。\lecturetimestamp{01:11:45--01:12:13}}
\end{classroomqa}

湮灭光子并没有消失。它们属于最终馈入微波背景的热辐射历史。然而，辐射在今天对能量收支贡献甚微，因为
\begin{equation}
\rho_{\rm rad}\propto a^{-4},
\qquad
\rho_m\propto a^{-3}.
\end{equation}
光子数很大，但每个光子的能量都会红移；物质则保留其静质量能量。

\begin{classroomqa}
\audiencequestion{如果湮灭产生了如此多的光子，为什么辐射在今天动力学上可以忽略不计？}
\lecturerresponse{光子仍然存在，主要以微波背景的形式，但它们的能量密度多损失了一个标度因子的幂次。膨胀既稀释了它们的数量，又使每个光子红移。\lecturetimestamp{01:12:14--01:13:18}}
\end{classroomqa}

\begin{classroomqa}
\audiencequestion{极重的媒介粒子是否意味着标准模型不完备？}
\lecturerresponse{大统一质子衰变的媒介粒子确实会是超出最小标准模型的新物理。后来关于标准模型重子反常破坏的更正并不能提供这样一种普通的质子衰变粒子，其本身也不能提供一个完整的定量重子生成模型。\lecturetimestamp{01:13:20--01:13:31}}
\end{classroomqa}

\begin{classroomqa}
\audiencequestion{两个示意性衰变通道之间若存在巨大差异，能否把最初纯\(p\)-\(\bar p\)粒子群转化为观测到的相等电子数与质子数？}
\lecturerresponse{人们可以编造这样的记账方式，但热平衡等离子体本就包含所有可及粒子与反粒子的大量粒子群。此外，观测到的\(CP\)不对称并不能成为让某一通道任意占主导的许可。一个自洽的模型必须演化完整的耦合反应网络，而不只是黑板上的两支箭头。\lecturetimestamp{01:13:33--01:15:10}}
\end{classroomqa}

重子生成讨论到此便诚实地结束了。定性的要素已知；观测到的数值尚未被推导出来。更完整的答案有待更完整的粒子理论以及对极早期宇宙的可控历史描述。

\section{第二个数：原初成团性}

暴胀还处理了另一个通过宇宙微波背景而变得定量化的问题：为什么宇宙在大尺度上如此均匀且各向同性？宇宙微波背景具有极其精确的黑体谱，温度接近\(2.7\,\mathrm K\)，而且该温度在每个方向上几乎相同。曾经只是一个宽泛的宇宙学印象，如今已成为高精度的事实。

引力使这种平滑性更加令人惊讶。扩散会把一滴墨水散开，直到其浓度变得均匀。吸引性的引力却相反：一个过密区域会把周围物质吸引过来，使其密度增加，同时使邻近区域变空。因此，微小的扰动会随时间增长，从而成为星系、星系团和空洞的种子。

\begin{figure}[htbp]
\centering
\begin{coursefitfigure}
\begin{tikzpicture}[x=1cm,y=1cm,>=Latex]
\draw[thick] (-4.2,0) ellipse (1.25 and 0.72);
\draw[very thick] (0,0) ellipse (1.55 and 0.92);
\draw[thick] (4.2,0) ellipse (1.25 and 0.72);
\node at (-4.2,0) {underdense};
\node at (0,0) {overdense};
\node at (4.2,0) {underdense};
\draw[->,thick] (-2.65,0)--(-1.65,0);
\draw[->,thick] (2.65,0)--(1.65,0);
\draw[->,thick] (-2.85,0.75)--(-1.45,0.4);
\draw[->,thick] (2.85,0.75)--(1.45,0.4);
\draw[->,thick] (-2.85,-0.75)--(-1.45,-0.4);
\draw[->,thick] (2.85,-0.75)--(1.45,-0.4);
\end{tikzpicture}
\end{coursefitfigure}
\caption{引力不稳定性把物质从低密度的周围区域转移到过密区域，从而放大初始反差，而不是将其扩散抹平。}
\label{fig:c08-gravity-growth}
\end{figure}

把这种增长倒推回去，意味着早期宇宙要平滑得多。完全均匀同样不可接受：完全对称会一直保持下去，任何星系都无法形成。有用的量度是密度反差
\begin{equation}
\delta\equiv\frac{\delta\rho}{\rho},
\end{equation}
其中\(\delta\rho\)是相对于背景密度的局域偏离。与Jim Peebles及其他早期宇宙学家相关的结构增长估计要求，在解耦时的原初振幅约为
\begin{equation}
\boxed{\frac{\delta\rho}{\rho}\sim10^{-5}}
\end{equation}
其中\(10^{-4}\)--\(10^{-5}\)被引作历史上的范围。这个数值再次使问题变得尖锐，却仍未解释其取值。

由Alan Guth在1980年提出的广义暴胀方案，是一段巨大的、近似指数式的膨胀时期。把一个小而平滑的区域拉伸许多个数量级，会稀释先前存在的曲率，并把不规则性推到大于可观测区域的尺度上。皱气球类比抓住了这种平滑化的直觉，但不应把它当作空间的实际嵌入。

\begin{classroomqa}
\audiencequestion{当附近物质在不同方向上明显成团时，宇宙怎么可能是各向同性的？}
\lecturerresponse{各向同性是一个大尺度统计陈述，并不是说这个房间、银河系或每条视线看起来都相同。地球局部多山，但与其半径相比极其平滑。同样，星系和星系团是叠加在非常平滑的大尺度背景上的局部结构。\lecturetimestamp{01:25:37--01:26:22}}
\end{classroomqa}

\begin{classroomqa}
\audiencequestion{随机过程难道不会偶尔造出极其过密的方向吗？}
\lecturerresponse{这些扰动并非不受限制的独立涨落。它们构成一个随机场，具有特定的、近似均匀且各向同性的统计谱。即使其波长很大，振幅也很小，量级为\(10^{-5}\)。该谱及其起源是尚待展开的主题，而不是任意大的抛硬币式偏离的假设。\lecturetimestamp{01:26:26--01:28:23}}
\end{classroomqa}

\begin{classroomqa}
\audiencequestion{暴胀是由早期对应的高温引起的吗？}
\lecturerresponse{温度并不是暴胀期间的主要变量。核心现象是由场的应力能量驱动的快速指数膨胀。当暴胀结束、其能量转化为粒子时，热浴才再次变得重要。\lecturetimestamp{01:28:25--01:29:04}}
\end{classroomqa}

\section{力学前奏：带阻力的运动}

为了给暴胀方程做准备，考虑一个单位质量的石块在黏性流体中下落。把它的竖直坐标记作\(\phi\)，以预示即将取代石块的标量场。若其势能为\(V(\phi)\)，则保守力为
\begin{equation}
F(\phi)=-\frac{dV}{d\phi},
\label{eq:c08-force-potential}
\end{equation}
而没有阻力时，牛顿方程为\(\ddot\phi=F\)。

\begin{figure}[htbp]
\centering
\includegraphics[width=0.82\textwidth]{lecture_08_figure_06.png}
\caption{落石类比中可读的受力关系 \(F=-dV/d\phi\)。向下的箭头确定了约定：\(\phi\) 增大时指向上方。}
\label{fig:c08-force-board}
\lectureframe{01:35:09.513}
\end{figure}

对于线性黏滞阻力，附加力为 \(-\gamma\dot\phi\)，其中 \(\gamma>0\)。方程变为
\begin{equation}
\ddot\phi+\gamma\dot\phi=F(\phi)
=-V'(\phi).
\label{eq:c08-viscous-motion}
\end{equation}
阻力在静止时为零，并且总是与速度方向相反。如果 \(\phi\) 向上增大，则下落的石块有 \(F<0\) 和 \(\dot\phi<0\)；于是 \(-\gamma\dot\phi>0\)，因此阻力指向上方，符合要求。

\begin{classroomqa}
\classquestion{为什么势力和阻力带有相同的负号？它们难道不应该指向相反方向吗？}
\lecturerresponse{它们的方向由不同的量决定。势力为 \(-V'(\phi)\)。阻力为 \(-\gamma\dot\phi\)。在向下运动过程中，\(F\) 和 \(\dot\phi\) 都为负，因此 \(F\) 指向下方，而 \(-\gamma\dot\phi\) 指向上方。所以代数运算在不改变任一公式的情况下给出了相反的物理方向。\lecturetimestamp{01:34:16--01:35:27}}
\end{classroomqa}

随着石块加速，阻力增大，直到各力平衡。对于近似恒定的 \(F\)，终端运动满足
\begin{equation}
0=F-\gamma\dot\phi_{\rm term},
\qquad
\dot\phi_{\rm term}=\frac{F}{\gamma}.
\end{equation}
符号是正确的：向下的 \(F\) 给出负的终端速度。因此，浅势阱和大的 \(\gamma\) 会产生沿山坡向下的缓慢运动。

\begin{figure}[htbp]
\centering
\begin{tikzpicture}[x=1cm,y=1cm,>=Latex]
\draw[->] (0,0)--(7.0,0) node[right] {\(\phi\)};
\draw[->] (0,0)--(0,4.4) node[above] {\(V(\phi)\)};
\draw[very thick] plot[smooth] coordinates {(0.5,3.7) (1.6,3.55) (2.8,3.2) (4.0,2.55) (5.1,1.65) (6.3,0.6)};
\filldraw (1.45,3.58) circle (0.09);
\draw[->,thick] (1.55,3.48)--(2.25,3.23);
\node[align=left,anchor=west] at (2.45,3.75) {large drag,\\slow downhill motion};
\end{tikzpicture}
\caption{浅势阱与强阻力。场的类比将用空间均匀的标量振幅取代石块的坐标。}
\label{fig:c08-shallow-potential}
\end{figure}

\section{暴胀子与哈勃摩擦}

暴胀假说引入一个额外的标量场 \(\phi\)，习惯上称为暴胀子。历史上，这是为解释宇宙的平滑性而提出的大胆假设。我们从近似空间均匀的场开始，\(\phi=\phi(t)\)。膨胀会拉伸空间变化，而均匀问题揭示了本质动力学。

对于正则标量场，完整的局域能量密度为
\begin{equation}
\rho_\phi
=
\frac12\dot\phi^{\,2}
+\frac{1}{2a^2}(\boldsymbol\nabla\phi)^2
+V(\phi).
\end{equation}
在均匀近似 \(\boldsymbol\nabla\phi=0\) 下，它简化为
\begin{equation}
\rho_\phi=\frac12\dot\phi^{\,2}+V(\phi).
\label{eq:c08-scalar-density}
\end{equation}
第一项是场位形空间中的动能：它衡量的是时间变化，而不是在普通空间中的平移运动。

\begin{classroomqa}
\audiencequestion{\(\phi\) 上的点表示对时间的偏导数吗？}
\lecturerresponse{是的。对于均匀场，\(\dot\phi=d\phi/dt\)。空间梯度会贡献一个单独的项，但在当前假设下它们为零。\lecturetimestamp{01:41:08--01:41:34}}
\end{classroomqa}

追踪一个单位共动坐标盒。它的物理体积为 \(a^3(t)\)，因此其中的标量能量为
\begin{equation}
E_\phi
=a^3(t)
\left[
\frac12\dot\phi^{\,2}+V(\phi)
\right].
\end{equation}
显式因子 \(a^3(t)\) 是新要素：即使共动坐标固定，物理盒子也会膨胀。

\begin{classroomqa}
\audiencequestion{标度因子 \(a\) 依赖于空间位置吗？}
\lecturerresponse{在均匀 FRW 模型中不依赖。在给定的宇宙时，它在整个空间中相同，但可以随时间变化。如果 \(\phi\) 依赖于位置，就必须恢复其空间梯度能量。\lecturetimestamp{01:44:12--01:44:49}}
\end{classroomqa}

\begin{classroomqa}
\audiencequestion{尽管我们现在只写一个 \(\phi(t)\)，场值仍然与空间中的每一点相关联吗？}
\lecturerresponse{是的。场为每个时空点赋予一个值。均匀性意味着这些空间值在每个时刻都相同，从而把动力学约化为一个共同的函数 \(\phi(t)\)。\lecturetimestamp{01:44:50--01:45:06}}
\end{classroomqa}

\begin{classroomqa}
\audiencequestion{场的动能项描述的是什么物理运动？场是在变强还是变弱？}
\lecturerresponse{这由运动方程、初始值和初始速度以及 \(V'(\phi)\) 的符号决定。类比是一个粒子，仅凭其能量公式无法知道它的运动方向。\lecturetimestamp{01:45:07--01:46:08}}
\end{classroomqa}

均匀拉格朗日量为动能减势能，
\begin{equation}
L_\phi=a^3(t)
\left[
\frac12\dot\phi^{\,2}-V(\phi)
\right].
\label{eq:c08-scalar-lagrangian}
\end{equation}
应用欧拉--拉格朗日方程得到
\begin{align}
\frac{d}{dt}\left(\frac{\partial L_\phi}{\partial\dot\phi}\right)
&=
\frac{\partial L_\phi}{\partial\phi},\\
\frac{d}{dt}\left(a^3\dot\phi\right)
&=-a^3V'(\phi).
\end{align}
乘积法则给出两项，
\begin{equation}
a^3\ddot\phi+3a^2\dot a\,\dot\phi
=-a^3V'(\phi).
\end{equation}
除以 \(a^3\) 并定义 \(H=\dot a/a\)，我们得到
\begin{equation}
\boxed{
\ddot\phi+3H\dot\phi+V'(\phi)=0
}
\label{eq:c08-hubble-friction}
\end{equation}
或者，用 \(F=-V'(\phi)\) 表示，
\begin{equation}
\ddot\phi+3H\dot\phi=F.
\end{equation}
每一行都来自方程~\eqref{eq:c08-scalar-lagrangian}中的膨胀体积因子。

\begin{classroomqa}
\audiencequestion{\(H=\dot a/a\) 在各处都相同吗？它真的是常数吗？}
\lecturerresponse{在均匀模型中，它在宇宙时切片上的每个空间点都相同。它不必随时间保持恒定；``哈勃常数''通常指其当前值，而 \(H(t)\) 是整个演化过程中的哈勃速率。\lecturetimestamp{01:50:32--01:50:57}}
\end{classroomqa}

方程~\eqref{eq:c08-hubble-friction}使这一类比变得精确。它与黏性运动具有相同的形式，其中阻尼系数为
\begin{equation}
\gamma_{\rm cosmic}=3H.
\end{equation}
因此，膨胀会减缓均匀标量场的演化。如果\(H\)很大，且\(V(\phi)\)足够平坦，则惯性项与阻尼项相比变得很小，场便趋近于慢滚关系
\begin{equation}
3H\dot\phi\simeq-V'(\phi).
\end{equation}
就像小球在冷而黏稠的油中下落一样，即使有力作用在场之上，场也可能需要极其漫长的时间才能穿过一个浅势阱。

下一步是把标量能量密度代入Friedmann方程，并让膨胀率对缓慢演化的场作出响应。这个耦合系统会产生暴胀。就目前而言，必要的准备已经完成：场储存动能和势能，膨胀的体积产生Hubble摩擦，而缓慢滚动的场能够使其势能几乎保持恒定足够长的时间，从而驱动快速膨胀。

\section{需要继续推进的内容}

重子不对称性和原初密度反差扮演着平行的角色。仅仅给\(10^{-8}\)或\(10^{-5}\)命名，并不能解释它们，但二者各自都把一种模糊的宇宙学印象转化为一个定量问题。前者导向Sakharov条件；后者则导向暴胀。

通向暴胀的桥梁同样精确。FRW宇宙中的均匀标量场满足
\[
\ddot\phi+3H\dot\phi+V'(\phi)=0.
\]
中间项并不是人为插入的唯象黏性。它来自物理体积\(a^3(t)\)的增长。仅凭这一结果，就能组织起下一讲中关于一个由沿其势能缓慢滚动的场驱动的宇宙的叙述。
